\section*{Bilan du projet}

Ce projet de fin d'études nous a permis de concevoir et de développer
\textbf{TAC-TIC IntelliPM}, une plateforme de gestion de projets complète et
intelligente. En cinq sprints Scrum de deux semaines chacun, nous avons livré
l'intégralité des 28 user stories planifiées, représentant 177 story points.

\medskip
Les principaux résultats techniques obtenus sont les suivants :

\begin{itemize}
  \item \textbf{Architecture distribuée opérationnelle} : backend Laravel~12
        exposant 251 routes REST, SPA React~19 intégralement en français, et
        microservice FastAPI Python~3.11 pour les modèles ML.
  \item \textbf{Sécurité multicouche} : authentification Sanctum (Bearer Token),
        RBAC à deux niveaux (système + projet), Policies Laravel, journal d'audit
        complet. Protection OWASP Top~10 couverte.
  \item \textbf{Temps réel} : notifications in-app et mises à jour Kanban via
        Laravel Reverb, latence médiane inférieure à 20\,ms.
  \item \textbf{IA opérationnelle} : AUC ROC 0,92 (détection de risques),
        MAE 1,42\,h (estimation de durée), inférence ML $<$ 45\,ms (P50).
  \item \textbf{Qualité certifiée} : SonarCloud Quality Gate passé --- 0 bug
        critique, 0 vulnérabilité, duplication $\approx$ 2\,\%.
  \item \textbf{Performance} : P95 API $<$ 500\,ms, temps de chargement initial
        SPA $<$ 2\,s (code splitting + cache React Query).
\end{itemize}

\section*{Déploiement}

L'application est conçue pour un déploiement cloud en trois composants indépendants :

\begin{table}[H]
  \centering
  \begin{tabularx}{\textwidth}{L{3cm}L{4cm}X}
    \toprule
    \textbf{Composant} & \textbf{Infrastructure} & \textbf{Détails} \\
    \midrule
    \rowcolor{lightgray}
    Frontend React  & Vercel (CDN mondial)    & Build Vite, déploiement automatique GitHub Actions \\
    Backend Laravel & VPS 4 vCPU / 8\,Go RAM & PHP-FPM + Nginx, Redis, Laravel Horizon (queues) \\
    \rowcolor{lightgray}
    ML FastAPI      & VPS 2 vCPU / 4\,Go RAM & Uvicorn, modèles persistés Joblib, rechargement à chaud \\
    Base de données & OVH MySQL 8 managé      & Sauvegardes automatiques, réplication lecture \\
    \bottomrule
  \end{tabularx}
\end{table}

La chaîne CI/CD (GitHub Actions) orchestre : lint, tests PHPUnit, analyse SonarCloud,
build Docker, déploiement automatique sur branche \texttt{main}. Cloudflare assure le
TLS, le CDN et la protection DDoS.

\section*{Apports personnels}

Ce projet nous a offert une formation transversale couvrant plusieurs domaines :

\begin{description}[leftmargin=3cm]
  \item[Génie logiciel] Architecture en couches, patterns Repositories et Services,
        SOLID, Clean Architecture.
  \item[Bases de données] Modélisation ERD, optimisation des requêtes Eloquent,
        migrations, indexation.
  \item[Intelligence artificielle] Préparation de données, feature engineering,
        sélection de modèles, validation croisée.
  \item[Sécurité] RBAC, politiques d'autorisation, protection OWASP, audit trail.
  \item[Méthodes agiles] Cérémonies Scrum, story points, vélocité, rétrospective.
  \item[DevOps] CI/CD GitHub Actions, Docker Compose, déploiement cloud.
\end{description}

\section*{Limites}

Nous identifions les principales limites de la version actuelle :

\begin{itemize}
  \item Les données d'entraînement des modèles ML sont synthétiques (simulées) ;
        en production réelle, l'historique de projets alimenterait des modèles
        plus précis.
  \item L'application est uniquement disponible en version web ; une application
        mobile (React Native) permettrait un accès en mobilité.
  \item La fonctionnalité d'intégration GitHub/GitLab (lier des commits à des tâches)
        a été reportée faute de temps.
\end{itemize}

\section*{Perspectives}

Les évolutions envisagées pour TAC-TIC IntelliPM s'articulent autour de trois axes :

\paragraph{Enrichissement fonctionnel}
Application mobile React Native, tableau de bord exécutif, chat intégré aux tâches,
localisation arabe (RTL), gamification (badges, points de productivité), export
PDF/Excel des rapports et du Gantt.

\paragraph{Intelligence artificielle avancée}
Réentraînement continu sur les données réelles, analyse NLP des commentaires pour
détection de conflits, valeurs SHAP pour l'explicabilité des recommandations,
détection automatique de surcharge d'équipe.

\paragraph{Industrialisation SaaS}
Déploiement Kubernetes/Helm, architecture multi-tenant, connecteurs Teams/Slack,
certification ISO~27001, modèle de facturation par siège.

\bigskip
\noindent\rule{\linewidth}{0.5pt}

\medskip
Ce projet a été l'occasion de cultiver l'équilibre entre rigueur technique
et humilité intellectuelle, entre précision de l'ingénieur et curiosité du
chercheur. Nous espérons que TAC-TIC IntelliPM constituera une base solide
pour une suite industrielle au service des équipes de développement tunisiennes
et au-delà.
